\section{The Craft of Dying Well: \textit{Ars Moriendi}}

``To die well is to die gladly and wilfully.''\endnote{\emph{The Book of the Craft of Dying, and Other Early English Tracts Concerning Death}, ed.\ Frances M.\ M.\ Comper (London: Longmans, Green, and Co., 1917), ch.\ 1, pp.\ 5--9 (\emph{Bene mori, est libenter mori}).} That crisp maxim opens one of the most influential pastoral manuals of late-medieval Christendom: the fifteenth-century \emph{Ars Moriendi}, rendered into Middle English as \emph{The Book of the Craft of Dying}. Drawing on Jean Gerson, Chancellor of Paris, together with St.\ Anselm, St.\ Gregory the Great, St.\ Bernard, and St.\ Augustine, it treated dying neither as biological failure nor as solitary mood, but as a craft to be learned in health and practiced with the support of the Christian community around the sickbed.

\subsection{The Commendation of Holy Death}

Measuring bodily death against the far graver catastrophe of spiritual death in mortal sin, the treatise sets Psalm 33:22 (\emph{Mors peccatorum pessima}, ``the death of the wicked is very evil'') beside Psalm 115:15 (\emph{Preciosa est in conspectu Domini mors sanctorum ejus}). Crucially, that preciousness belongs not to martyrs alone, but even to lifelong sinners ``if they die in the state of very repentance and contrition, and in the very faith, and virtue, and charity of Holy Church'' (pp.\ 5--6). For those who die in the Lord, death is ``a going out of prison, and an ending of exile; a discharging of an heavy burden, that is the body; finishing of all infirmities; a scaping of all perils\ldots{} paying of the debt of natural duty; turning again into his country; and entering into bliss and joy'' (p.\ 6). To know how to die---in Middle English, ``to con die''---is ``to have an heart and a soul every ready up to Godward, that when-that-ever death come, he may be found all ready; withouten any retraction receive him, as a man would receive his well-beloved and trusty friend and fellow, that he had long abideth and looked after'' (pp.\ 8--9). Every Christian must therefore live so as to ``have his life in patience, and his death in desire'' (p.\ 9, citing Phil.\ 1:23).

\subsection{The Five Temptations of the Dying}

At the center of the \emph{Ars Moriendi} stands a sober map of the final spiritual combat: ``men that die, in their last sickness and end, have greatest and most grievous temptations, and such as they never had before in all their life.''\endnote{\emph{Book of the Craft of Dying}, ch.\ 2, pp.\ 9--21.} Five assaults are principal:

\begin{enumerate}
\item \emph{Temptation against Faith.} Because faith is the foundation of salvation (Heb.\ 11:6), the adversary first tries to unsettle belief with doubt, superstition, or heresy. The remedy is public and direct: bystanders recite the Creed aloud ``with an high voice oft times'' to steady the sick person in the Church's unity and drive away demons (pp.\ 10--11).

\item \emph{Desperation (Despair).} Weighting bodily agony with remembrance of unconfessed or grievous sins, the tempter seeks ``to superadd sorrow to sorrow.'' Even if one had committed as many sins ``as be drops of water in the sea, and gravel stones in the strand,'' and lacked time or speech for confession, true contrition with the will to confess saves him (Ps.\ 50:19). Beside Peter, Paul, Mary Magdalene, and the Good Thief, the treatise sets St.\ Bernard's contemplation of Christ on the Cross:

\begin{quote}
Take heed and see: His head is inclined to salve thee; His mouth to kiss thee; His arms spread to be-clip thee; His hands thrilled to give thee; His side opened to love thee; His body along strait to give all Himself to thee. (p.\ 14)
\end{quote}

In the Middle English wording, ``be-clip'' means embrace and ``thrilled'' means pierced.

\item \emph{Impatience.} Acute pain and long illness tempt the unprepared heart to murmur against God and wound charity. Against impatience stand Christ's promise, ``In your patience ye shall possess your souls'' (Luke 21:19), St.\ Gregory's maxim that God sends temporal chastisement here lest He inflict everlasting vengeance hereafter, St.\ Augustine's prayer ``Here cut, here burn, so that Thou spare me everlastingly,'' and St.\ Albert the Great's counsel that final illness borne willingly serves as a purgatory before death (pp.\ 15--18).

\item \emph{Complacence (Spiritual Pride).} Failing to break a devout soul by unbelief, despair, or impatience, the enemy turns to flattery: ``O how stable art thou in the faith! how strong in hope! how sad in patience! O how many good deeds hast thou done!'' Quoting St.\ Gregory---that pride over good deeds casts a man down before the Author of humility---the manual prescribes St.\ Antony's double poise: when pushed upward by pride, hold yourself down by meekness; when dragged downward by despair, hold yourself up by hope (pp.\ 18--19).

\item \emph{Attachment to Temporal Things.} Disordered solicitude for property, riches, or even earthly ties to spouse, children, and friends pulls the heart backward. Letting go of outward cares and consenting willingly to the pain of death in obedience to God---the manual cites Duns Scotus---satisfies for venial sins and remits temporal punishment (pp.\ 19--20).
\end{enumerate}

Over all five temptations stands the liberating anchor of 1 Corinthians 10:13: the devil ``may compel no man, nor in no manner of wise prevail against him for to consent to him---as long as a man hath the use of reason with him---but if he will wilfully consent unto him'' (pp.\ 20--21).

\subsection{The Interrogations, the Cross, and the Warning to the Healthy}

Chapter 3 supplies bedside interrogations attributed to St.\ Anselm and Jean Gerson to be asked while speech and reason endure.\endnote{\emph{Book of the Craft of Dying}, ch.\ 3, pp.\ 22--27; ch.\ 4, pp.\ 27--31.} Gerson's seven questions examine Catholic faith and obedience, honest confession of sin, contrition from love of God, firm purpose of amendment, forgiveness of all enemies, restitution of ill-gotten goods, and total trust in Christ's Passion (pp.\ 24--26). St.\ Anselm's interrogation rises to one of medieval piety's most fearless cries:

\begin{quote}
Lord, I put the death of Our Lord Jesu Christ between me and mine evil deeds, between me and the judgment; otherwise will I not strive with Thee. If He say: Thou hast deserved damnation; say thou again: The death of our Lord Jesu Christ I put between me and mine evil merits, and the merits of His worthy passion I offer for merits I should have had, and alas I have not. (p.\ 23)
\end{quote}

Then the dying person---or the brethren around the bed---repeats three times \emph{In manus tuas, Domine}. Chapter 4 adds that the dying should imitate Christ's five deeds on the Cross: praying, crying from the heart, weeping in repentance, commending His soul to the Father, and willingly surrendering His spirit (pp.\ 27--28).

Chapter 5 directs its sharpest realism at the living and those who tend the sick. It denounces the cruel flattery of hiding mortal danger under false promises of recovery---citing Gerson's warning that ``feigned behoting of bodily heal'' sends unprepared souls to ruin, whereas Isaias spoke plainly to sick King Ezechias---and begs every reader to learn the craft while well:\endnote{\emph{Book of the Craft of Dying}, ch.\ 5, pp.\ 32--39; ch.\ 6, pp.\ 39--47.}

\begin{quote}
For sooth, dear sister or brother, I tell thee sooth, believe me thereof, that when death or great sickness falleth upon thee, devotion passeth out from thee; and the more near they take thee and grip thee, the further fleeth devotion from thee\ldots{} Therefore if thou wilt not be deceived or err---if thou wilt be sure---do busily what thou mayst while thou art in heal, and hast the use and freedom of thy five wits and reason well disposed, and while thou mayst be master of thyself and of thy deeds. (pp.\ 37--38)
\end{quote}

When a brother enters his agony, the treatise praises the monastic custom of striking the wooden board so that every member drops his work and hastens to the bedside with litanies, the Crucifix, and holy water (pp.\ 36, 38--39).

\subsection{St.\ Robert Bellarmine and the School of the Psalms}

More than a century later, St.\ Robert Bellarmine gathered this entire tradition into his late treatise \emph{De arte bene moriendi} (1620), governed by one enduring rule: he who lives well dies well, and no one may safely postpone that craft to the hour of final weakness. The scriptural roots of Bellarmine's wisdom shine throughout his \emph{Commentary on the Book of Psalms}.\endnote{St.\ Robert Bellarmine, \emph{A Commentary on the Book of Psalms}, trans.\ John O'Sullivan (Dublin: James Duffy, 1866), on Ps.\ 38:4--6, 12--13, and Ps.\ 89:4--15.}

Expounding David's prayer in Psalm 38---``O Lord, make me know my end, and what is the number of my days''---Bellarmine explains that the Psalmist seeks not idle foreknowledge of a calendar date, but grace to pierce the illusion of earthly permanence:

\begin{quote}
Man walks and passes through life in the image, not in the reality of things, having before him on his journey, not the realities, but the images and the shadows. This life is but an image of the happy life that alone is the true one; the health of this life is only an image of the immortality that alone deserves the name of health; the beauty of this world is only the shadow of the beauty with which we will be clothed when ``the just shall shine like the sun in the kingdom of their father.'' (on Ps.\ 38:6)
\end{quote}

Confessing himself ``a stranger with thee and a sojourner'' (Ps.\ 38:12), David begs ``O forgive me, that I may be refreshed before I go hence''---because, Bellarmine notes, ``if you do not forgive me here, I will not go to rest, but to prison.'' In Psalm 89, measuring our seventy or eighty years---passing like morning grass and spun out ``as a spider''---against God's eternity, Bellarmine shows how even the infirmities of age humble youthful pride and train the soul under God's hand for the morning when His servants are filled with mercy and bless even the years in which their patience was proved.
